\section*{Bab \ref{ch:b2:powerseries} --- Deret Pangkat}

\begin{solution}{exo:b2:powerseries:1}
$\sum \frac{n^2}{2^n}z^n$: rasionya $\frac{(n+1)^2}{2^{n+1}}\cdot
\frac{2^n}{n^2} \to \frac12$, jadi $R = 2$.

$\sum \frac{z^n}{\binom{2n}{n}}$: karena $\binom{2n}{n} \sim
\frac{4^n}{\sqrt{\pi n}}$ (\cref{ex:b2:comparison:centralbinomial}),
maka $\abs{a_n}^{-1} \approx 4^n$ hingga faktor polinomial: jadi $R = 4$
(lewat uji rasio: $\frac{\binom{2n}{n}}{\binom{2n+2}{n+1}} =
\frac{(n+1)^2}{(2n+1)(2n+2)} \to \frac14$).

$\sum z^{n!}$: koefisiennya $a_k = 1$ bila $k = n!$, dan $0$ bila tidak.
Untuk $\abs z < 1$, $\sum \abs z^{n!}$ konvergen (didominasi geometri);
sedangkan untuk $\abs z \geq 1$ sukunya tak menuju $0$: jadi $R = 1$.

$\sum (2 + (-1)^n)^n z^n$: koefisiennya $3^n$ (untuk $n$ genap) dan $1$
(untuk $n$ ganjil). Keterbatasan $a_nr^n$ menuntut $3r \leq 1$; dan $r <
\frac13$ berhasil: jadi $R = \frac13$.
\end{solution}

\begin{solution}{exo:b2:powerseries:2}
Ketiganya berjari-jari $1$ (lewat uji rasio). Di $z = 1$: $\sum 1$
divergen; $\sum\frac1n$ divergen; $\sum\frac{1}{n^2}$ konvergen. Di
$z = -1$: $\sum(-1)^n$ divergen; $\sum\frac{(-1)^n}{n}$ konvergen
(berselang-seling); $\sum\frac{(-1)^n}{n^2}$ konvergen (\omterm{def:b2:series:def}{mutlak}).
Jadi kelakuan di perbatasan tak terlihat oleh jari-jarinya.
\end{solution}

\begin{solution}{exo:b2:powerseries:3}
Dari $\frac{1}{1-x} = \sum x^n$, turunkan lalu kalikan dengan $x$
(\cref{thm:b2:powerseries:calculus}):
\[
\sum n x^n = \frac{x}{(1-x)^2} .
\]
Turunkan sekali lagi, lalu kalikan lagi dengan $x$:
\[
\sum n^2 x^n = x\,\frac{\dd}{\dd x}\Bigl(\frac{x}{(1-x)^2}\Bigr)
= \frac{x(1 + x)}{(1-x)^3} .
\]
Jumlah ketiganya: ia bagian ganjil dari $-\ln(1 - x)$:
\[
\sum_{n\geq0} \frac{x^{2n+1}}{2n+1}
= \frac{-\ln(1-x) + \ln(1+x)}{2}
= \frac12 \ln\frac{1+x}{1-x}
= \operatorname{artanh} x .
\]
\end{solution}

\begin{solution}{exo:b2:powerseries:4}
Pecahan parsialnya: $\frac{1}{(1-x)(2-x)} = \frac{1}{1-x} -
\frac{1}{2 - x} = \sum x^n - \frac12\sum \bigl(\frac x2\bigr)^n$, jadi
\[
\frac{1}{(1-x)(2-x)} = \sum_{n\geq0} \Bigl(1 -
\frac{1}{2^{n+1}}\Bigr)x^n,
\qquad R = 1 .
\]

$\ln(1 + x + x^2) = \ln\frac{1 - x^3}{1 - x} = \ln(1 - x^3) -
\ln(1 - x) = \sum_{n\geq1}\frac{x^n}{n} -
\sum_{m\geq1}\frac{x^{3m}}{m}$: jadi koefisien $x^n$ adalah
$\frac1n$ bila $3 \nmid n$, dan $\frac1n - \frac{3}{n} = -\frac2n$ bila
$3 \mid n$. Jari-jarinya $1$ (yakni halangan terdekatnya: deret bagi
$\ln(1-x^3)$).
\end{solution}

\begin{solution}{exo:b2:powerseries:5}
\omterm{thm:b2:series:fubini}{Hasil kali Cauchy} $\sum_{m \geq 0} x^m$ (yang koefisiennya $1$) dan
$\sum_{k\geq1} \frac{x^k}{k}$ (yang koefisiennya $\frac1k$ untuk $k \geq
1$), dan keduanya konvergen \omterm{def:b2:series:def}{mutlak} untuk $\abs x < 1$: maka koefisien
$x^n$ pada hasil kalinya adalah $\sum_{k=1}^{n} \frac1k \cdot 1 = H_n$.
Karena itu
\[
\Bigl(\sum x^m\Bigr)\Bigl(\sum \frac{x^k}{k}\Bigr)
= \frac{1}{1-x}\cdot\bigl(-\ln(1-x)\bigr)
= \sum_{n\geq1} H_n x^n .
\]
\end{solution}

\begin{solution}{exo:b2:powerseries:6}
Kalikan rekurensinya dengan $x^{n+1}$ lalu jumlahkan (dengan $\abs x <
\frac12$):
\[
U(x) - 1 = 2x\,U(x) + \sum_{n\geq0} n\,x^{n+1}
= 2x\,U(x) + \frac{x^2}{(1-x)^2} ,
\]
lewat \cref{exo:b2:powerseries:3}. Karena itu
\[
U(x) = \frac{1}{1 - 2x}\Bigl(1 + \frac{x^2}{(1-x)^2}\Bigr)
= \frac{1 - 2x + 2x^2}{(1-2x)(1-x)^2} .
\]
Pecahan parsialnya (menutupi di $x = \frac12$ memberi koefisien $2$;
sedangkan di kutub gandanya $x = 1$ koefisiennya $-1$; dan
koefisien tengahnya lenyap setelah dinilai di $x = 0$):
\[
U(x) = \frac{2}{1-2x} - \frac{1}{(1 - x)^2} .
\]
Setelah keduanya diuraikan:
\[
u_n = 2\cdot 2^n - (n + 1) = 2^{n+1} - n - 1 .
\]
(Periksa: $u_0 = 1$, $u_1 = 2u_0 + 0 = 2 = 4 - 2$.)
\end{solution}

\begin{solution}{exo:b2:powerseries:7}
Secara induktif: $f'(x) = \frac{2}{x^3}\eu^{-1/x^2}$, dan jika $f^{(n)}(x)
= P_n(\tfrac1x)\eu^{-1/x^2}$ maka
\[
f^{(n+1)}(x) =
\Bigl(-\frac{1}{x^2}\,P_n'\Bigl(\frac1x\Bigr) +
\frac{2}{x^3}\,P_n\Bigl(\frac1x\Bigr)\Bigr)\eu^{-1/x^2} :
\]
yang kembali berbentuk seperti yang dinyatakan. Di $0$: hasil bagi selisihnya $\frac{f^{(n)}(h)}{h} =
\frac1h P_n(\frac1h)\eu^{-1/h^2} \to 0$ ketika $h \to 0$, sebab
$Q(u)\,\eu^{-u^2} \to 0$ ketika $u \to \pm\infty$ untuk sembarang polinomial
$Q$ (eksponensial mengalahkan pangkat): jadi secara induktif setiap $f^{(n)}(0)$ ada
dan lenyap, lalu tiap $f^{(n)}$ \omterm{def:b2:metric:continuity}{kontinu} di $0$ lewat limit yang
sama. Jadi $f \in C^\infty$ dengan deret Taylor nol di $0$; dan
deret Taylornya berjumlah $0 \neq f$: sehingga tak \omterm{def:b2:powerseries:analytic}{analitik} di $0$.
\end{solution}

\begin{solution}{exo:b2:powerseries:8}
Deret binomialnya: $\sqrt{1-4x} = \sum_{k\geq0}
\binom{1/2}{k}(-4x)^k$. Untuk $k = n + 1 \geq 1$:
\begin{align*}
\binom{1/2}{n+1}(-4)^{n+1}
&= \frac{\frac12\bigl(\frac12 - 1\bigr)\cdots\bigl(\frac12 -
n\bigr)}{(n+1)!}\,(-4)^{n+1}\\
&= \frac{(-1)^n\,1\cdot3\cdots(2n-1)}{2^{n+1}(n+1)!}\,(-4)^{n+1}
= -\frac{2}{n+1}\cdot\frac{(2n)!}{n!\,n!} ,
\end{align*}
memakai $1\cdot3\cdots(2n-1) = \frac{(2n)!}{2^n n!}$. Karena itu
\[
C(x) = \frac{1 - \sqrt{1-4x}}{2x}
= \frac{1}{2x}\sum_{n\geq0}\frac{2}{n+1}\binom{2n}{n}x^{n+1}
= \sum_{n\geq0} \frac{1}{n+1}\binom{2n}{n}\,x^n :
\]
$C_n = \frac{1}{n+1}\binom{2n}{n}$. Jari-jarinya $\frac14$ (yakni deret
binomial dalam $4x$). Asimtotiknya lewat
\cref{ex:b2:comparison:centralbinomial}:
\[
C_n \sim \frac{4^n}{\sqrt{\pi}\; n^{3/2}} .
\]
\end{solution}

\begin{solution}{exo:b2:powerseries:9}
Dengan $A_n = \sum_{k \leq n} a_k \to A$: \omterm{thm:b2:series:abel}{penjumlahan Abel} memberi, untuk
$0 \leq x < 1$,
\[
\sum_{n=0}^{\infty} a_n x^n = (1 - x)\sum_{n=0}^{\infty} A_n x^n
\]
(kedua ruasnya konvergen sebab $(A_n)$ terbatas; dan kesamaannya menyusul
dari $a_n = A_n - A_{n-1}$ beserta pengindeksan ulang). Karena $(1 -
x)\sum x^n = 1$:
\[
\sum_n a_nx^n - A = (1-x)\sum_{n} (A_n - A)x^n .
\]
Diberikan $\varepsilon$, pilihlah $N$ dengan $\abs{A_n - A} \leq
\varepsilon$ untuk $n > N$; maka
\[
\Bigl|\sum a_nx^n - A\Bigr|
\leq (1-x)\sum_{n \leq N}\abs{A_n - A} + \varepsilon(1 -
x)\sum_{n > N}x^n
\leq (1-x)\,C_N + \varepsilon ,
\]
dan setelah dilewatkan $x \to 1^-$: limsupnya $\leq \varepsilon$ untuk
setiap $\varepsilon$. Jadi limit radialnya $A$.

Penerapannya: $\sum \frac{(-1)^{n-1}}{n}$ konvergen (berselang-seling),
dan untuk $x < 1$ deret pangkatnya berjumlah $\ln(1 + x)$: jadi menurut
Abel jumlahnya $\ln 2$. Demikian pula $\sum\frac{(-1)^n}{2n+1}x^{2n+1} =
\arctan x$ memberi $\frac\pi4$ di $x = 1$ --- jadi bukti integral
Tahun ke-1 kini menjadi struktural.
\end{solution}

\begin{solution}{exo:b2:powerseries:10}
Baik $\sum\frac{x^n}{n}$ maupun $\sum\frac{x^n}{n+1}$ berjari-jari
$1$, dan $\frac{1}{n(n+1)} = \frac1n - \frac1{n+1}$, jadi untuk $0
< \abs x < 1$:
\[
\sum_{n\geq1}\frac{x^n}{n(n+1)}
= -\ln(1-x) - \frac1x\sum_{n\geq1}\frac{x^{n+1}}{n+1}
= -\ln(1-x) - \frac{-\ln(1-x) - x}{x}
= 1 + \frac{1-x}{x}\ln(1-x) .
\]
Jari-jarinya $1$; dan $\norm{x^n/(n(n+1))}_{\infty,\intcc{-1}1} =
\frac{1}{n(n+1)}$ \omterm{def:b2:series:summable}{terjumlahkan}: jadi kekonvergenannya normal pada
$\intcc{-1}{1}$, sehingga jumlahnya \omterm{def:b2:metric:continuity}{kontinu} di sana. Ketika $x \to
1^-$, berlaku $(1-x)\ln(1-x) \to 0$ dan bentuk tertutupnya menuju $1$ ---
sepakat dengan nilai teleskopnya $\sum\frac{1}{n(n+1)} =
\lim_N\bigl(1 - \frac{1}{N+1}\bigr) = 1$ di $x = 1$.
\end{solution}

\begin{solution}{exo:b2:powerseries:11}
Menyortir $n!$ permutasinya menurut himpunan titik tetapnya: memilih
$k$ titik tetapnya (ada $\binom nk$ cara) lalu mengacaukan $n - k$
objek sisanya memberi $n! = \sum_{k=0}^n\binom nk D_{n-k}$.
Setelah dibagi $n!$:
\[
1 = \sum_{k=0}^{n}\frac{1}{k!}\cdot\frac{D_{n-k}}{(n-k)!} ,
\]
yang persis mengatakan bahwa hasil kali Cauchy $\eu^x =
\sum\frac{x^k}{k!}$ dan $D(x) = \sum D_n\frac{x^n}{n!}$ adalah
$\sum x^n = \frac{1}{1-x}$. Kedua faktornya konvergen \omterm{def:b2:series:def}{mutlak}
untuk $\abs x < 1$ (sebab $D_n \leq n!$, jadi $D$ didominasi
deret geometri): maka kesamaan hasil kalinya sah
(\cref{prop:b2:powerseries:operations}), dan
\[
D(x) = \frac{\eu^{-x}}{1-x} .
\]
\omterm{thm:b2:series:fubini}{Hasil kali Cauchy} $\eu^{-x} = \sum\frac{(-1)^kx^k}{k!}$ dan
$\sum x^m$: koefisien $x^n$-nya adalah
$\sum_{k=0}^{n}\frac{(-1)^k}{k!}$, dan berkat ketunggalan koefisien
deret pangkat (\cref{thm:b2:powerseries:calculus}):
\[
\frac{D_n}{n!} = \sum_{k=0}^{n}\frac{(-1)^k}{k!}
\xrightarrow[n\to\infty]{} \eu^{-1} :
\]
jadi sekitar $37\%$ dari semua permutasi tak punya titik tetap, berapa
pun $n$.
\end{solution}

\begin{solution}{exo:b2:powerseries:12}
Deret binomial dengan $\alpha = -\frac12$ di $-4x$:
\[
\binom{-1/2}{n}(-4)^n
= \frac{\bigl(-\frac12\bigr)\bigl(-\frac32\bigr)\cdots
\bigl(-\frac{2n-1}2\bigr)}{n!}(-4)^n
= \frac{1\cdot3\cdots(2n-1)}{2^n\,n!}\,4^n
= \frac{(2n)!}{2^n n!}\cdot\frac{2^n}{n!}
= \binom{2n}{n},
\]
memakai $1\cdot3\cdots(2n-1) = \frac{(2n)!}{2^nn!}$. Karena itu
$(1-4x)^{-1/2} = \sum\binom{2n}nx^n$ untuk $\abs{4x} < 1$.
Setelah dikuadratkan (lewat hasil kali Cauchy, yang sah berkat
kekonvergenan \omterm{def:b2:series:def}{mutlaknya}) lalu dibandingkan dengan
$\frac{1}{1-4x} = \sum 4^nx^n$: koefisien $x^n$ pada
kuadratnya adalah $\sum_{k=0}^n\binom{2k}k\binom{2n-2k}{n-k}$, sehingga
ketunggalan koefisiennya memberi
\[
\sum_{k=0}^{n}\binom{2k}{k}\binom{2n-2k}{n-k} = 4^n .
\]
\end{solution}

\begin{solution}{pb:b2:powerseries:1}
\textbf{1.} Dengan $a_n = r_n - r_{n+1}$, penjumlahan parsialnya:
\[
\sum_{n=N}^{M} a_nx^n
= r_Nx^N + \sum_{n=N+1}^{M} r_n\bigl(x^n - x^{n-1}\bigr)
- r_{M+1}x^M .
\]
Untuk $0 \leq x \leq 1$ pertambahan $x^{n-1} - x^n$ bersifat
taknegatif dan berteleskop menjadi $x^N - x^M$; jadi dengan $s = \sup_{n\geq
N}\abs{r_n}$:
\[
\Bigl|\sum_{n=N}^{M}a_nx^n\Bigr|
\leq s\bigl(x^N + (x^N - x^M) + x^M\bigr) = 2s\,x^N \leq 2s .
\]

\textbf{2.} Karena $r_n \to 0$, berlaku $\sup_{n\geq N}\abs{r_n} \to 0$:
jadi pertanyaan 1 persis merupakan kriteria Cauchy \omterm{def:b2:funcseq:def}{seragam} pada
$\intcc{0}{1}$, sehingga $\sum a_nx^n$ konvergen \omterm{def:b2:funcseq:def}{seragam} di sana dan
jumlahnya \omterm{def:b2:metric:continuity}{kontinu} (\cref{thm:b2:funcseq:seriestransfer}).
Karena nilainya di $1$ adalah $\sum a_n$, kekontinuan di $1$ tak lain
limit radial pada \cref{exo:b2:powerseries:9}.

\textbf{3.} Untuk $\abs x < 1$ ketiga deret pangkatnya konvergen
\omterm{def:b2:series:def}{mutlak} dan $\bigl(\sum a_nx^n\bigr)\bigl(\sum b_nx^n\bigr)
= \sum c_nx^n$ (\cref{prop:b2:powerseries:operations}). Menurut
pertanyaan 2, tiap faktornya dan ruas hasil kalinya \omterm{def:b2:metric:continuity}{kontinu} pada
$\intcc{0}{1}$ (sebab deret koefisiennya konvergen menurut
hipotesisnya); lalu melewatkan $x \to 1^-$ pada kesamaannya: $AB = C$.

\textbf{4.} Di sini
\[
\abs{c_n} = \sum_{k=0}^{n}
\frac{1}{\sqrt{(k+1)(n-k+1)}}
\geq \sum_{k=0}^{n}\frac{2}{n+2}
= \frac{2(n+1)}{n+2} \geq 1,
\]
lewat AM--GM: $\sqrt{(k+1)(n-k+1)} \leq \frac{(k+1) +
(n-k+1)}{2} = \frac{n+2}{2}$. Jadi suku umum $\sum c_n$
tidak menuju $0$: sehingga \omterm{thm:b2:series:fubini}{hasil kali Cauchynya} divergen, padahal
kedua faktornya konvergen (yakni deret berselang-seling).

\textbf{5.} \cref{exo:b2:powerseries:5} memberi
$\frac{-\ln(1-x)}{1-x} = \sum H_nx^n$ (untuk $\abs x < 1$).
Antiturunan suku demi sukunya
(\cref{thm:b2:powerseries:calculus}~(2)), dengan kedua ruasnya lenyap
di $0$, memberi
\[
\frac{\bigl(\ln(1-x)\bigr)^2}{2}
= \sum_{n\geq1}\frac{H_n}{n+1}\,x^{n+1} .
\]
Ketururunannya: $(n+2)H_n \geq (n+1)H_{n+1}$ setara dengan $H_n \geq
1$, yang benar untuk $n \geq 1$; dan $\frac{H_n}{n+1} \sim \frac{\ln
n}{n} \to 0$: jadi di $x = -1$ deretnya konvergen menurut
uji berselang-seling. Setelah disulih $x \mapsto -x$ lalu diterapkan
pertanyaan 2 di $x = 1$:
\[
\frac{(\ln 2)^2}{2}
= \sum_{n\geq1}\frac{H_n}{n+1}(-1)^{n+1},
\]
yakni nilai yang diumumkan.

\textbf{6.} $f(x) = \sum(-1)^nx^n = \frac{1}{1+x} \to \frac12$
ketika $x \to 1^-$: jadi terjumlahkan-Abel ke $\frac12$. Tetapi jumlah
parsialnya $1, 0, 1, 0, \dots$: sehingga divergen.

\textbf{7.} Diberikan $\varepsilon > 0$, pilihlah $m$ dengan $\abs{u_n}
\leq \varepsilon$ untuk $n > m$; lalu untuk $N \geq m$:
\[
\Bigl|\frac{u_1 + \dots + u_N}{N}\Bigr|
\leq \frac{\abs{u_1} + \dots + \abs{u_m}}{N} +
\varepsilon\,\frac{N - m}{N}
\leq \frac{C_m}{N} + \varepsilon,
\]
jadi $\limsup \leq \varepsilon$ untuk setiap $\varepsilon$: sehingga
rata-ratanya menuju $0$.

\textbf{8.} Untuk $0 \leq x \leq 1$: $1 - x^n = (1-x)(1 + x +
\dots + x^{n-1}) \leq n(1-x)$, jadi
\[
\Bigl|\sum_{n=0}^N a_n(1 - x_N^n)\Bigr|
\leq (1 - x_N)\sum_{n=1}^N n\abs{a_n}
= \frac1N\sum_{n=1}^{N}n\abs{a_n} .
\]
Untuk $n > N$: $\abs{a_n} = \frac{n\abs{a_n}}{n} \leq
\frac{1}{N}\sup_{m>N}m\abs{a_m}$, dan $\sum_{n>N}x_N^n \leq
\frac{1}{1 - x_N} = N$, sehingga
\[
\Bigl|\sum_{n>N}a_nx_N^n\Bigr|
\leq \frac{\sup_{m>N}m\abs{a_m}}{N}\cdot N
= \sup_{m>N}\,m\abs{a_m} .
\]

\textbf{9.} Uraikanlah
\[
A_N - L = \sum_{n=0}^{N}a_n\bigl(1 - x_N^n\bigr)
- \sum_{n>N}a_nx_N^n + \bigl(f(x_N) - L\bigr) .
\]
Suku pertamanya menuju $0$ menurut pertanyaan 7 (yakni rata-rata
$n\abs{a_n} \to 0$), yang kedua menurut pertanyaan 8 (sebab supnya menuju
$0$), dan yang ketiga karena $x_N \to 1^-$ dan $f(x) \to L$. Jadi
$A_N \to L$: itulah teorema Tauber.

\textbf{10.} Untuk $x \in \intco{0}{1}$ dan sembarang $N$ berlaku
$\sum_{n\leq N}a_nx^n \leq f(x) \leq M$ (sebab sukunya taknegatif).
Lewatkan $x \to 1^-$ pada jumlah berhingganya: $\sum_{n\leq N}a_n \leq M$.
Jadi jumlah parsialnya naik dan terbatas: sehingga $\sum a_n$
konvergen, lalu pertanyaan 2 memberi $\lim_{x\to1^-}f(x) =
\sum a_n$.

\textbf{11.} Sebab $\sigma_N - L$ adalah rata-rata dari $N$ bilangan
$A_n - L$ (dengan $0 \leq n < N$), yang menuju $0$: yakni pertanyaan 7.

\textbf{12.} $A_n = 1$ untuk $n$ genap dan $0$ untuk yang ganjil, jadi
$A_0 + \dots + A_{N-1} = \lceil N/2\rceil$, sehingga $\sigma_N =
\frac{\lceil N/2\rceil}{N} \to \frac12$, yakni nilai Abel pada
pertanyaan 6.

\textbf{13.} Di bawah $\sigma_N \to L$ diperoleh $S_n = O(n)$, sehingga
$A_n = S_n - S_{n-1} = O(n)$ dan $a_n = O(n)$: jadi semua deret
di bawah berjari-jari $\geq 1$. Untuk $\abs x < 1$, dari $a_n = A_n -
A_{n-1}$ dan $A_nx^n \to 0$:
\[
(1-x)\sum_n A_nx^n = \sum_n A_nx^n - \sum_n A_nx^{n+1}
= \sum_n a_nx^n = f(x),
\]
dan secara identik $(1-x)\sum S_nx^n = \sum A_nx^n$, jadi $f(x) =
(1-x)^2\sum_n S_nx^n = (1-x)^2\sum_n(n+1)\sigma_{n+1}x^n$.
Akhirnya $\sum(n+1)x^n = \frac{1}{(1-x)^2}$
(\cref{exo:b2:powerseries:3}), dan itulah kesamaan keduanya.

\textbf{14.} Setelah $L$ kali kesamaan keduanya dikurangkan dari yang
pertama:
\[
f(x) - L = (1-x)^2\sum_{n\geq0}(n+1)
\bigl(\sigma_{n+1} - L\bigr)x^n .
\]
Diberikan $\varepsilon$, pilihlah $N$ dengan $\abs{\sigma_{n+1} - L}
\leq \varepsilon$ untuk $n \geq N$; maka
\[
\abs{f(x) - L} \leq (1-x)^2 C_N +
\varepsilon(1-x)^2\sum_{n}(n+1)x^n
= (1-x)^2C_N + \varepsilon ,
\]
lalu setelah dilewatkan $x \to 1^-$: $\limsup \leq \varepsilon$. Jadi
$f(x) \to L$: itulah teorema Frobenius.

\textbf{15.} $f(x) = \sum(-1)^n(n+1)x^n = \frac{1}{(1+x)^2}$
(turunkan deret geometrinya di $-x$): jadi nilai Abelnya
$\frac14$. Jumlah parsialnya: $A_{2k} = k+1$ dan $A_{2k+1} = -(k+1)$
(lewat induksi seketika). Lalu $S_{2m-1} = 0$ (sebab pasangan
berurutannya saling meniadakan) dan $S_{2m} = m + 1$, jadi
\[
\sigma_{2m} = \frac{S_{2m-1}}{2m} = 0,
\qquad
\sigma_{2m+1} = \frac{m+1}{2m+1} \to \frac12 :
\]
jadi $(\sigma_N)$ punya dua nilai gugus yang berbeda: sehingga tidak
terjumlahkan-Cesàro. Bersama pertanyaan 6 dan 11--14, hierarki
konvergen $\Rightarrow$ Cesàro $\Rightarrow$ Abel bersifat tegas
pada kedua panahnya.

\textbf{16.} Deret antiturunannya $F(x) =
\sum\frac{a_n}{n+1}x^{n+1}$ berjari-jari sama dan $F' = f$ pada
$\intco{0}{1}$ (\cref{thm:b2:powerseries:calculus}); dan karena
$\sum\frac{a_n}{n+1}$ konvergen, Bagian I (pertanyaan 2) membuat $F$
\omterm{def:b2:metric:continuity}{kontinu} pada $\intcc{0}{1}$. Karena $\int_0^x f = F(x)$ (sebab
turunannya sama dan nilainya sama-sama $0$ di $0$), berlaku
\[
\int_0^x f \xrightarrow[x\to1^-]{} F(1)
= \sum_{n\geq0}\frac{a_n}{n+1} :
\]
jadi \omterm{def:b2:integration:improper}{integral tak wajarnya} ada dengan nilai yang dinyatakan.

\textbf{17.} $\frac{\ln(1+x)}{x} =
\sum_{n\geq1}\frac{(-1)^{n-1}}{n}x^{n-1}$ (berjari-jari $1$ dan
\omterm{def:b2:metric:continuity}{kontinu} di $0$). Deret bagi $\frac{a_m}{m+1}$ adalah
$\sum_{n\geq1}\frac{(-1)^{n-1}}{n^2}$, yang konvergen \omterm{def:b2:series:def}{mutlak}: jadi
pertanyaan 16 memberi $\int_0^1\frac{\ln(1+x)}{x}\dd x = \eta$. Lalu di
dalam $\sum\frac1{n^2}$ yang konvergen \omterm{def:b2:series:def}{mutlak}, kelompokkan ulang yang
genap dan yang ganjil:
\[
\eta = \sum_{\text{ganjil}}\frac1{n^2} -
\sum_{\text{genap}}\frac1{n^2}
= \sum_{n}\frac1{n^2} - 2\sum_{k}\frac1{(2k)^2}
= \Bigl(1 - \frac12\Bigr)\sum_n\frac1{n^2}
= \frac12\sum_{n\geq1}\frac1{n^2} .
\]

\textbf{18.} Menurut Leibniz: $\frac{1}{3n+1}\downarrow0$, jadi deretnya
konvergen. Pada $\intco{0}{1}$ berlaku $\sum(-1)^nx^{3n} =
\frac{1}{1+x^3}$, dan $\sum\frac{(-1)^n}{3n+1}$ konvergen: jadi
pertanyaan 16 memberi $\sum\frac{(-1)^n}{3n+1} =
\int_0^1\frac{\dd x}{1+x^3}$. Pecahan parsialnya (periksa:
$\frac13(x^2-x+1) + \frac{2-x}{3}(1+x) = 1$):
\[
\int_0^1\frac{\dd x}{1+x^3}
= \frac13\ln2 + \frac13\int_0^1\frac{2-x}{x^2-x+1}\dd x .
\]
Setelah menulis $2 - x = -\frac12(2x-1) + \frac32$: bagian
$\ln(x^2-x+1)$-nya lenyap di kedua ujungnya, dan
\[
\frac32\int_0^1\frac{\dd x}{(x-\frac12)^2 + \frac34}
= \frac32\cdot\frac{2}{\sqrt3}
\Bigl[\arctan\frac{2x-1}{\sqrt3}\Bigr]_0^1
= \sqrt3\cdot\frac{\pi}{3} = \frac{\pi}{\sqrt3} .
\]
Totalnya: $\frac13\bigl(\ln2 + \frac{\pi}{\sqrt3}\bigr)$.

\textbf{19.} Setelah menyulihkan $t = x^2$ pada deret
\cref{exo:b2:powerseries:12} lalu mengintegralkannya suku demi
suku (sebab antiturunan $(1-x^2)^{-1/2}$ yang lenyap di $0$ adalah
$\arcsin$):
\[
\arcsin x = \sum_{n\geq0}
\frac{\binom{2n}n}{4^n(2n+1)}x^{2n+1}
\qquad(\abs x < 1) .
\]
Koefisiennya $\sim \frac{1}{2\sqrt\pi\,n^{3/2}}$
(\cref{ex:b2:comparison:centralbinomial}), yang \omterm{def:b2:series:summable}{terjumlahkan}: jadi
deretnya konvergen \emph{normal} pada $\intcc{-1}{1}$, jumlahnya
\omterm{def:b2:metric:continuity}{kontinu} di sana, lalu berimpit dengan $\arcsin$ yang \omterm{def:b2:metric:continuity}{kontinu}
pada $\intoo{-1}{1}$, sehingga juga di $x = 1$:
\[
\sum_{n\geq0}\frac{\binom{2n}n}{4^n(2n+1)}
= \arcsin 1 = \frac\pi2 .
\]

\textbf{20.} $C_n4^{-n} \sim \frac{1}{\sqrt\pi\,n^{3/2}}$
(\cref{exo:b2:powerseries:8}): jadi \omterm{def:b2:funcseq:series}{kekonvergenan normal} $\sum
C_nx^n$ pada $\intcc{0}{\frac14}$, sehingga jumlahnya \omterm{def:b2:metric:continuity}{kontinu}
di sana; lalu pada $\intoo{0}{\frac14}$ ia sama dengan
$\frac{1-\sqrt{1-4x}}{2x}$
(\cref{ex:b2:powerseries:catalan}), yang berlimit di
$\frac14^-$ sebesar $\frac{1-0}{1/2} = 2$. Jadi $\sum_{n\geq0}
C_n4^{-n} = 2$.

\textbf{21.} $f(x) = \sum_{n\geq m}a_nx^n = x^m g(x)$ dengan
$g(x) = \sum_{k\geq0}a_{m+k}x^k$; dan jika $(a_nr^n)$ terbatas maka
$(a_{m+k}r^k)$ juga (bagilah dengan $r^m$): jadi $g$ berjari-jari $\geq
R$. Lalu $g$ \omterm{def:b2:metric:continuity}{kontinu} dengan $g(0) = a_m \neq 0$, jadi $g \neq 0$
pada suatu $\intcc{-\delta}{\delta}$, sehingga $f(x) = x^mg(x) \neq 0$
untuk $0 < \abs x \leq \delta$.

\textbf{22.} Fungsi $d = f - h$ adalah jumlah deret pangkat di dekat $0$
yang lenyap di titik tak nol $x_k \to 0$. Andaikan suatu
koefisien $d$ tak nol, maka pertanyaan 21 akan memberi
persekitaran $0$ yang tertusuk dan bebas dari nol milik $d$ ---
yang bertentangan dengan $d(x_k) = 0$. Jadi semua koefisien $d$ lenyap:
sehingga $f$ dan $h$ berkoefisien sama dan berimpit di dekat $0$.

\textbf{23.} Fungsi $h(x) = \frac{1}{1+x^2} =
\sum(-1)^nx^{2n}$ (berjari-jari $1$) memenuhi $h(\frac1k) =
\frac{1}{1 + 1/k^2} = \frac{k^2}{k^2+1}$. Sembarang $f$ \omterm{def:b2:powerseries:analytic}{analitik} dengan
nilai yang sama sepakat dengan $h$ di titik $\frac1k \to 0$:
jadi menurut teorema identitas (pertanyaan 22), $f = \frac{1}{1+x^2}$
di dekat $0$ --- yakni penyelesaian tunggalnya.

\textbf{24.} Misalkan $Z$ himpunan titik $I$ yang punya
persekitaran tempat $f$ lenyap secara identik: ia \omterm{def:b2:metric:topology}{terbuka} menurut
definisinya, dan tak kosong (yakni subintervalnya). Lalu tertutup di $I$:
sebab jika $y \in I$ limit titik-titik $Z$, maka $y$ titik akumulasi
nol milik $f$; dan setelah menguraikan $f$ menjadi deret pangkat di $y$
(berkat keanalitikannya) lalu menerapkan pertanyaan 21--22 yang dipusatkan
ulang di $y$, semua koefisien di $y$ lenyap, jadi $f \equiv 0$ di dekat $y$:
sehingga $y \in Z$. Karena interval bersifat \omterm{def:b2:metric:connected}{terhubung}, maka $Z = I$:
jadi $f \equiv 0$ pada $I$. Khususnya, \omterm{def:b2:powerseries:analytic}{fungsi analitik} pada $\R$ yang
lenyap di luar sebuah himpunan \omterm{def:b2:metric:compact}{kompak} lenyap pada sebuah interval,
sehingga lenyap di mana-mana: jadi tak ada bonggol \omterm{def:b2:powerseries:analytic}{analitik} yang tak nol.
Sedangkan dunia $C^\infty$ berbeda: menempelkan fungsi datar pada
\cref{exo:b2:powerseries:7} (misalnya $x \mapsto
\eu^{-1/x^2}\mathbf 1_{x>0}$ beserta cerminannya) menghasilkan bonggol
mulus yang berpenyangga \omterm{def:b2:metric:compact}{kompak}.

\textbf{25.} (i) \omterm{def:b2:funcseq:series}{Kekonvergenan normal} hidup pada subcakram \omterm{def:b2:metric:compact}{kompak}
yang tegas di dalam cakramnya; sedangkan teorema Abel memperluas
kekontinuannya ke sebuah titik perbatasan, dengan satu-satunya hipotesis
bahwa deret koefisiennya konvergen di sana. (ii) Konversnya berlaku
di bawah syarat Tauber $na_n \to 0$ (pertanyaan 9), dan
keterjumlahan Cesàro duduk tegas di antara kekonvergenan dan
keterjumlahan Abel (yakni Frobenius, pertanyaan 14--15). (iii) Untuk
seorang teman: $\sum\frac{(-1)^n}{3n+1} = \int_0^1\frac{\dd x}{1+x^3}$
lewat mengintegralkan deret geometrinya sampai ke perbatasan, lalu
pecahan parsial. (iv) Rata-rata Cesàro kembali pada bab Fourier
sebagai teorema Fejér, tempat merata-ratakan jumlah parsialnya
memperbaiki kegagalan \omterm{def:b2:funcseq:def}{kekonvergenan titik demi titik} --- jadi obat yang
sama, pasien yang baru.

\end{solution}
